%----------Abstract (English)---------------------------------------------------
\section*{Abstract}
In autonomous driving research, synthetic environments are widely used to train and evaluate perception systems. However, discrepancies between map-generation processes can introduce a persistent domain gap between different synthetic environments. This thesis investigates the domain gap between two CARLA environments representing the same region in Ingolstadt: (i) an automatically generated 3D map derived from OpenStreetMap via an OSM$\rightarrow$OpenDRIVE$\rightarrow$CARLA pipeline, and (ii) a manually authored reference map.

The domain gap is analyzed primarily through its \textbf{structural} component (road geometry, topology, lane connectivity, and junction structure). To ensure experimental integrity, the OSM extraction bounds are fixed, the full toolchain is versioned, and pipeline determinism is evaluated using an external audit. The audit reveals non-determinism at the byte level while preserving the audited topological counts ($CV=0.0$ for road and junction counts), motivating a bounded analysis of observed structural variability under fixed inputs rather than any claim of intentional domain randomization.

Structural comparison between the auto-generated and manually authored Ingolstadt maps yields a whole-network geometry RMSE of 54.84\,m, a matched-subset RMSE of 24.47\,m (73\% of sampled road points within a 50\,m criterion), and a Hausdorff distance of 1{,}663\,m after CRS-normalized rigid alignment. The auto map is 4.88$\times$ longer in total road length and has 6.55$\times$ more junctions. The curvature distribution of the auto map is substantially wider than the manual reference (standard deviation 0.865 vs.\ 0.055\,m$^{-1}$), indicating high-curvature outlier segments absent in the hand-authored map. The corrected curvature KL divergence is $0.00031$ (dominated by the concentration of both distributions near $kappa pprox 0$); the decisive signal is therefore the standard-deviation gap, not the near-zero KL value. An earlier measurement of KL~$= 2{,}493.53$ was an artefact of a sampler that ignored 	exttt{paramPoly3} geometry elements; the corrected paramPoly3-aware value is reported throughout. A lane and junction connectivity analysis finds that the auto map carries no road-level predecessor/successor link declarations, relying entirely on junction-based connectivity, while all declared links in the manual map are fully resolvable.

Perception evidence is bounded rather than complete: the governed strict-evidence attempt on the auto-generated OpenDRIVE world targeted 50 frames but did not yield a passing final evidence pack, while a separate Town10 baseline confirmed rig-callback readiness on a built-in map. Perception capture on the manually-authored cooked Grid0828 map was bounded by CARLA runtime instability during live sensor spawn on the already-loaded custom world. All identified repo-local control-flow defects were diagnosed and corrected; nevertheless, the final governed attempt entered the rig-attach path with a valid client but produced zero frames, with streaming transport refusing connections persistently. This is classified as a CARLA runtime boundary rather than a remaining code defect. Perceptual domain gap quantification and generalization experiments therefore remain future work. An orthogonal semantic-object gap is observed in the authoritative contract-run artifact: the four-category object counter records zero barriers, buildings, poles, and trees in the auto map versus 1{,}538, 992, 3{,}427, and 4{,}232 respectively in the manual reference. A subsequent pipeline fix confirmed that building injection now produces the expected objects; that result is retained as supplementary pipeline evidence and does not alter the authoritative measurement.
